%webamc CODE=r203/2023/tcp
%webamc TITLE=TCP
%webamc RND=0

On considère dans cet exercice un processus $C$ qui s'adresse
successivement à deux serveurs TCP $S_1$ et $S_2$. Le premier lui
refuse la connexion, et le second accepte. Dans son échange avec
$S_2$, $C$ envoie une requête de 100 octets à laquelle $S_2$ répond
par 1470 octets de données. $S_2$ met ensuite fin à la connexion.

Voici les chronogrammes retraçant les échanges entre $C$ et $S_1$ (à
gauche) et entre $C$ et $S_2$ (à droite).

\begin{minipage}[t]{.40\textwidth}
  \begin{center}~\\\begin{tempDiag}[>=stealth'][numbers=0]{2}{5cm}
  \tempDiagNode{1}{3pt}{C}{$C$}
  \tempDiagNode{2}{3pt}{S}{$S_1$}
  %
  \tempDiagMsg{1}{2}{}\node[left of=pt-1,number] {A};
  \tempDiagMsg{2}{1}{}\node[right of=pt-2,number] {B};
\end{tempDiag}
\end{center}
\end{minipage}
\hfill
\begin{minipage}[t]{.40\textwidth}
  \begin{center}~\\\begin{tempDiag}[>=stealth']{2}{6cm}
  \tempDiagNode{1}{3pt}{C}{$C$}
  \tempDiagNode{2}{3pt}{S}{$S_2$}
  %
  \tempDiagMsg{1}{2}{Seq=1000}
  \tempDiagMsg{2}{1}{Seq=200, SeqA=1001}
  \tempDiagMsg{1}{2}{SeqA=201}
  \tempDiagMsg{1}{2}{Len=100 (la requête)}
  \tempDiagMsg{2}{1}{}
  \tempDiagMsg{2}{1}{Len=??? (début de la réponse)}
  \tempDiagMsg{1}{2}{}
  \tempDiagMsg{2}{1}{Len=??? (fin de la réponse)}
  \tempDiagMsg{1}{2}{}
  \tempDiagMsg{2}{1}{}
  \tempDiagMsg{1}{2}{}
  \tempDiagMsg{2}{1}{}
\end{tempDiag}
\end{center}
\end{minipage}

Les paquets contenant les données (requête et réponse) sont les
paquets 4, 6 et 8. Les autres paquets servent uniquement à contrôler
l'échange (connexion, déconnexion, \ldots).  On rappelle que:
\begin{itemize}
\item {\em Seq} = numéro de séquence d'envoi, {\em SeqA} = numéro de
  séquence d'acquittement ;
\item en-tête IP = 20 octets, en-tête TCP = 20 octets, en-tête UDP = 8
  octets, et en-tête + FCS ethernet = 26 octets ;
\item une trame Ethernet ne peut pas encapsuler plus de 1~500 octets.
\end{itemize}

